\documentclass{article}
\usepackage{amsmath,amssymb,hyperref}
\begin{document}
\title{A complete pivotal obstruction for R6\_9\_8785}
\author{}\date{1 October 2026}\maketitle

\paragraph{Result and scope.}
The ring R6\_9\_8785 has no pivotal, spherical, pseudo-unitary, or unitary
categorification over $\mathbb C$. This is the pseudo-unitary ingredient
of the complete general exclusion in \texttt{general\_proof\_8785.tex}. The tensor digest is
\texttt{7897af45c48cb1d416710f37ea9c4c11226490421317d193c85e3cc5022b3da8}.
The accompanying verifier uses exact integer pairs and rational linear algebra;
it imports no census classification or induction algorithms.

\paragraph{Exact character table.}
Put $\phi=(1+\sqrt5)/2$, $t=8+13\phi=\phi^7$, and
$D=5(1+t^2)=1170+1885\phi$.
The six characters, in the input tensor's basis order, are
\[
\begin{array}{c|rrrrrr}
 &0&1&2&3&4&5\\\hline
 c_0&1&t&t-1&t-1&t+1&t+1\\
 c_1&1&21-13\phi&20-13\phi&20-13\phi&22-13\phi&22-13\phi\\
 c_2&1&1+2\phi&-1-\phi&1&\phi&-1-2\phi\\
 c_3&1&1+2\phi&1&-1-\phi&-1-2\phi&\phi\\
 c_4&1&3-2\phi&\phi-2&1&1-\phi&2\phi-3\\
 c_5&1&3-2\phi&1&\phi-2&2\phi-3&1-\phi
\end{array}
\]
The verifier checks all $6\cdot6^2$ multiplication equations. Six distinct
characters exhaust the rank-six algebra. The first is positive and hence is
FPdim. Their codegrees are $D,D',25+10\sqrt5,25+10\sqrt5,
25-10\sqrt5,25-10\sqrt5$, respectively. Consequently their central simple
dimensions $\beta_j=D/f_j$ in a hypothetical pseudo-unitary category are
\[
(1,\ 233+377\phi,\ 26+39\phi,\ 26+39\phi,
442+715\phi,\ 442+715\phi).
\]

\paragraph{Small-column anchoring criterion.}
For a commutative pseudo-unitary fusion category, label the summands $Z_j$ of
$I(1)$ by their characters $c_j$ and let $v_j$ be the restriction column of
$Z_j$. Then $v_{j0}=1$, $d(v_j)=\beta_j$, and
\[
 \beta_i c_i(v_j)=\beta_j c_j(v_i). \tag{1}
\]
Indeed, the character pulled back to the center is
$c_i(FW)=S_{Z_i,W}/\beta_i$, by the character-level correspondence preceding
Ostrik's Theorem 2.13, so (1) is symmetry of the modular $S$ matrix.
When small-dimensional columns have short exhaustive domains, fixing their
\emph{labeled} choices first makes (1) a rational linear system for every
other column. This is a general exact pruning method; no choice of labels
among equal codegrees is omitted.

Here $\beta_2=\beta_3=3t+2$. The coefficient of $\phi$ in
$d(v)=3t+2$, with $v_0=1$, forces $\sum_{i=1}^5v_i=3$.
The complete eight-column domain is
\[
\begin{split}
 &(1,0,0,1,0,2),\ (1,0,0,1,1,1),\ (1,0,0,1,2,0),\\
 &(1,0,1,0,0,2),\ (1,0,1,0,1,1),\ (1,0,1,0,2,0),\\
 &(1,2,0,0,0,1),\ (1,2,0,0,1,0).
\end{split}
\]
Exhaust all $8^2=64$ ordered labeled pairs $(v_2,v_3)$.
For each remaining column $v_j$, the equation $v_{j0}=1$, the two rational
components of its dimension equation, and the four rational components of
(1) with $i=2,3$ form a seven-by-six rational system of rank six.
There are precisely ten nonnegative integral completions. The verifier
also checks the complete row-sum and pairwise reciprocity identities for
every completion. Thus the ten frames are an exhaustive domain, not
numerically proposed induction data.

\paragraph{The twist-moment bound.}
For a nontrivial self-dual simple $X$ of dimension $d>0$, let $\ell$ be the
positive dimension mass in $I(X)$ of its summands also occurring in $I(1)$.
Their twists are one. The remaining positive weights have total $Dd-\ell$.
Ostrik's Theorems 2.5 and 2.7 (the latter from Ng--Schauenburg) give their
first real twist moment $-\ell$ and second real twist moment
$D\nu_2(X)-\ell$, where $\nu_2(X)=\pm1$. Since
$\cos^2\theta=(1+\cos2\theta)/2$, Cauchy--Schwarz gives
\[
 2\ell^2\le (Dd-\ell)\bigl(D(d+1)-2\ell\bigr).
\]
Cancellation yields the necessary inequality
\[
 (3d+1)\ell\le Dd(d+1). \tag{2}
\]
No choice of second-indicator sign is assumed: replacing it by $+1$ only
weakens the bound.

For $X=X_1$ ($d=t$), the ten frames give just three possible low masses:
\[
\begin{array}{c|c|c}
\ell&\text{number of frames}&(3t+1)\ell-Dt(t+1)\\\hline
14473+23413\phi&6&257812+417144\phi\\
13641+22061\phi&2&184284+298168\phi\\
12809+20709\phi&2&110756+179192\phi
\end{array}
\]
Every gap has two strictly positive integer coefficients. Thus every
possible frame contradicts (2), proving pseudo-unitary nonexistence.

\paragraph{Spherical and pivotal lifts.}
A spherical dimension function is one of the six characters. The four
characters $c_2,\ldots,c_5$ have global dimension of norm $125$. The FP
codegree has norm $21125=125\cdot169$. Every formal codegree divides the
global dimension (Ostrik Corollary 2.14); their quotient would have norm
$1/169$, so cannot be an algebraic integer in $\mathbb Q(\sqrt5)$.
The only remaining dimension characters are $c_0,c_1$, which are Galois
conjugate. A categorification on either branch therefore has a Galois
conjugate with all positive dimensions and would be pseudo-unitary,
already excluded. Finally all six tensor simples are self-dual. Every
pivotal structure is therefore spherical, proving the stated pivotal
obstruction. Unitary categorifications are pseudo-unitary and are excluded.

\paragraph{Primary sources.}
V. Ostrik, \emph{Pivotal fusion categories of rank 3}, arXiv:1309.4822,
Theorems 2.5, 2.7, 2.13 and Corollary 2.14,
\url{https://arxiv.org/pdf/1309.4822}.
S.-H. Ng and P. Schauenburg, \emph{Frobenius--Schur indicators and exponents
of spherical categories}, arXiv:math/0601012, Theorem 4.1,
\url{https://arxiv.org/pdf/math/0601012}.

\paragraph{Replay.}
Run \texttt{python -O verify\_8785.py}.
The machine-readable exact frames are in
\texttt{verified\_certificate.json}; the input is the sibling \texttt{tensor\_8785.json}. Neither numerical eigenvalues, LP solvers,
root-of-unity bounds, nor incomplete induction searches enter the proof.
\end{document}
